\documentclass[11pt]{article}

\usepackage[margin=1in]{geometry}
\usepackage{amsmath,amssymb,amsthm,mathtools}
\usepackage{microtype}
\usepackage[hidelinks]{hyperref}

\newtheorem{theorem}{Theorem}
\newtheorem{lemma}[theorem]{Lemma}
\newtheorem{proposition}[theorem]{Proposition}
\newtheorem{corollary}[theorem]{Corollary}
\theoremstyle{remark}
\newtheorem{remark}[theorem]{Remark}

\newcommand{\E}{\mathbb E}
\newcommand{\ind}{\mathbf 1}
\newcommand{\cU}{\mathcal U}
\newcommand{\cO}{\mathcal O}
\newcommand{\cV}{\mathcal V}
\newcommand{\Rstar}{\mathfrak R}
\DeclareMathOperator{\Unif}{Unif}

\title{Tight dynamic regret for trading profit with estimates of the suffix maximum}
\author{}
\date{}

\begin{document}
\maketitle

\begin{abstract}
At round $t$, the current price $p_t$ and an estimate $L_t$ of the suffix
maximum $M_t$ are shown before the player posts a quote.  A quote $A_t>p_t$
earns the net profit $A_t-p_t$ if a later price reaches $A_t$.  We measure
dynamic regret against the profit of the best quote on every round.  If the
total amounts by which the estimates are too low and too high are at most
$u$ and $o$, respectively, the tight worst-case rate is
\[
  \Theta\!\left(
    \min\left\{T-1,\,
      u+\min\bigl\{(T-1)o,\sqrt{(T-1)o}\bigr\}
    \right\}
  \right).
\]
The fixed-quote minimax regret has the same rate.  There is no logarithmic
term: exact estimates give zero dynamic regret on every sequence.  The term
caused by overestimation is linear when $(T-1)o$ is small and has the usual
square-root form once $(T-1)o$ is at least a constant.
\end{abstract}

\section{Model}

Fix a horizon $T\ge 1$ and put $N:=T-1$.  The environment fixes prices
$p_1,\ldots,p_T\in[0,1]$ and estimates $L_1,\ldots,L_N\in[0,1]$ before play.
For $t=1,\ldots,T$, let
\[
  M_t:=\max_{t\le s\le T}p_s,
\]
the suffix maximum of the price sequence.

At each round $t\le N$:
\begin{enumerate}
  \item the player is shown $p_t$ and an estimate $L_t$ of $M_t$;
  \item the player either abstains or posts a quote $A_t\in(p_t,1]$;
  \item the quote earns $A_t-p_t$ if $p_s\ge A_t$ for some $s>t$, and zero
        otherwise.
\end{enumerate}
Different rounds are treated as separate trades; there is no restriction on
the number of outstanding quotes.  The player may observe prices, fills, and
past rewards as they occur.  The terminal price $p_T$ is revealed only to
settle outstanding quotes; no new quote is posted at $T$.  An unfilled trade
contributes zero and is not liquidated at the horizon.  Abstention is explicit
because otherwise there would be no valid action when $p_t=1$.

Both $p_t$ and the reported estimate are therefore available before $A_t$
is chosen.  If the reported estimate is below $p_t$, the player replaces it
by $p_t$.  Since $p_t\le M_t$, this moves the estimate closer to $M_t$ and
cannot increase either underestimation or overestimation.  We reuse the
symbol $L_t$ for the value after this preprocessing.  Thus, without loss of
generality,
\begin{equation}\label{eq:estimate-normalization}
  p_t\le L_t\le1.
\end{equation}

Execution is non-strict: equality with a later price is enough.  A genuine
quote $a>p_t$ executes if and only if
\begin{equation}\label{eq:threshold-equivalence}
  a\le M_t.
\end{equation}
Indeed, if the suffix maximum is at least $a$, it cannot be attained at
round $t$ because $p_t<a$ and must be attained later.  We can therefore write
the round-$t$ reward as
\begin{equation}\label{eq:reward}
  r_t(a):=(a-p_t)\ind\{p_t<a\le M_t\},
\end{equation}
with reward zero for abstention.  The strict inequality $p_t<a$ and the
non-strict inequality $a\le M_t$ are both important.

\paragraph{Dynamic benchmark and dynamic regret.}
The largest profit attainable separately on round $t$ is $M_t-p_t$: quote
$M_t$ when $M_t>p_t$, and abstain when $M_t=p_t$.  The optimal dynamic
payoff is therefore
\begin{equation}\label{eq:dynamic-payoff}
  D_T(p):=\sum_{t=1}^{N}(M_t-p_t).
\end{equation}
This is the payoff of a clairvoyant player allowed to choose a different
quote on every round.  No policy can earn more: every successful quote is at
most $M_t$, and every failed quote earns zero.  For a randomized policy
$\pi$, its dynamic regret is
\begin{equation}\label{eq:dynamic-regret}
  R_T^{\mathrm{dyn},\pi}(p,L)
  :=D_T(p)-\E_\pi\sum_{t=1}^{N}r_t(A_t).
\end{equation}
This regret is nonnegative on every sequence.

\paragraph{Fixed-quote benchmark.}
For comparison, a fixed quote $a\in[0,1]$ is posted on rounds with $p_t<a$;
on the other rounds it means abstention.  Its best payoff is
\begin{equation}\label{eq:benchmark}
  B_T(p):=\max_{a\in[0,1]}
     \sum_{t=1}^{N}(a-p_t)\ind\{p_t<a\le M_t\}.
\end{equation}
This is a fixed absolute quote, not a fixed markup above $p_t$.  Since every
term earned by a fixed quote is at most $M_t-p_t$,
\begin{equation}\label{eq:benchmark-dominated}
  B_T(p)\le D_T(p).
\end{equation}
Define the corresponding fixed-quote regret by
\begin{equation}\label{eq:fixed-regret}
  R_T^{\mathrm{fix},\pi}(p,L)
  :=B_T(p)-\E_\pi\sum_{t=1}^{N}r_t(A_t).
\end{equation}
It can be negative because a changing policy can outperform every fixed
quote.  Equations~\eqref{eq:dynamic-regret}--\eqref{eq:fixed-regret} give the
pointwise comparison
\begin{equation}\label{eq:regret-order}
  R_T^{\mathrm{fix},\pi}(p,L)
  \le R_T^{\mathrm{dyn},\pi}(p,L).
\end{equation}

\paragraph{Directional errors.}
We use the same inclusive maximum $M_t$ as in the gross-revenue formulation:
\begin{equation}\label{eq:errors}
  \cU_T:=\sum_{t=1}^{N}(M_t-L_t)_+,
  \qquad
  \cO_T:=\sum_{t=1}^{N}(L_t-M_t)_+.
\end{equation}
Thus $\cU_T$ is total underestimation and $\cO_T$ is total overestimation.
The normalization $L_t\ge p_t$ also gives
\begin{equation}\label{eq:under-bounded-by-dynamic}
  \cU_T\le\sum_{t=1}^{N}(M_t-p_t)=D_T.
\end{equation}
For $b\in\{\mathrm{dyn},\mathrm{fix}\}$ and $0\le u,o\le N$, let
\begin{equation}\label{eq:minimax}
  \Rstar_T^b(u,o)
  :=\inf_\pi
      \sup_{\substack{p,L:\ p_t\le L_t,\
                       \cU_T\le u,\ \cO_T\le o}}
      R_T^{b,\pi}(p,L).
\end{equation}
The policies in the infimum may know $T,u,o$.  In fact, the policies below
use only $T$ and $o$.  The environment is oblivious: it fixes $(p,L)$ before
seeing the player's random choices, although $L_t$ may depend on future
prices.  By \eqref{eq:regret-order},
\begin{equation}\label{eq:minimax-order}
  \Rstar_T^{\mathrm{fix}}(u,o)
  \le \Rstar_T^{\mathrm{dyn}}(u,o).
\end{equation}
When the policy is clear from context, we omit $\pi$ from the regret
notation.

\section{Policies and upper bounds}

We use one of the following two policies.

\paragraph{Full subtraction.}
If the announced overestimation budget is $o$, propose
\begin{equation}\label{eq:full-subtraction}
  A_t=L_t-o
\end{equation}
when this value is above $p_t$; otherwise abstain.

\paragraph{Random subtraction.}
Fix $\gamma>0$.  Independently draw
$Z_t\sim\Unif[0,\gamma]$ and post
\begin{equation}\label{eq:random-subtraction}
  A_t=L_t-Z_t
\end{equation}
when this value is above $p_t$; otherwise abstain.

The full subtraction is best for very small $o$.  Random subtraction is
better once the accumulated cost of subtracting $o$ on every round becomes
large.

\begin{theorem}[Upper bounds]\label{thm:upper}
For every price and estimate sequence, the dynamic regret of random
subtraction satisfies, for every $\gamma>0$,
\begin{equation}\label{eq:random-upper}
  R_T^{\mathrm{dyn}}
  \le
  \cU_T+\frac{N\gamma}{2}+\frac{\cO_T}{\gamma}.
\end{equation}
On the class $\cO_T\le o$, full subtraction satisfies
\begin{equation}\label{eq:full-upper}
  R_T^{\mathrm{dyn}}\le \cU_T+No.
\end{equation}
Consequently, for $N\ge1$,
\begin{equation}\label{eq:minimax-upper}
  \Rstar_T^{\mathrm{dyn}}(u,o)
  \le
  \min\left\{
    N,\,
    u+\min\bigl\{No,\sqrt{2No}\bigr\}
  \right\}.
\end{equation}
When $o=0$, the full-subtraction policy uses no subtraction and satisfies
$R_T^{\mathrm{dyn}}\le\cU_T$.
\end{theorem}

\begin{proof}
We bound the shortfall from $D_T$ round by round.  We first prove
\eqref{eq:random-upper}.  Fix a round and abbreviate
$p=p_t$, $m=M_t$, $\ell=L_t$, and $Z=Z_t$.  The best round-by-round payoff
is $d:=m-p$.  Notice that $\ell\ge p$.

Suppose first that $\ell\le m$.  Whenever $\ell-Z>p$, the proposed quote is
no larger than $m$ and therefore executes.  Its reward is
$(\ell-Z-p)_+$, and
\begin{equation}\label{eq:under-round}
  d-(\ell-Z-p)_+
  \le (m-\ell)+Z.
\end{equation}

Now suppose that $\ell>m$ and put $\delta:=\ell-m$.  If $Z\ge\delta$, the
proposed quote is no larger than $m$.  If it is above $p$, it executes and the loss
against $d$ is $Z-\delta$; if it is at or below $p$, the player abstains and
$Z\ge \ell-p=\delta+d$, so the loss $d$ is again at most $Z-\delta$.  Thus the
loss on $\{Z\ge\delta\}$ is at most $Z$.  On $\{Z<\delta\}$ the proposed
quote exceeds $m$ and fails, but the loss is at most one.  If
$\delta\le\gamma$, this gives
\[
  \E[\text{loss at round }t]
  \le \E\bigl[Z\ind\{Z\ge\delta\}\bigr]+\Pr(Z<\delta)
  \le \frac\gamma2+\frac\delta\gamma.
\]
If $\delta>\gamma$, the loss is at most one, which is at most
$\delta/\gamma$.  Combining this with \eqref{eq:under-round},
\[
  d-\E r_t(A_t)
  \le
  (m-\ell)_++\frac\gamma2+\frac{(\ell-m)_+}{\gamma}.
\]
Summing over $t$ proves \eqref{eq:random-upper}.

For full subtraction, $(L_t-M_t)_+\le\cO_T\le o$ on every round.  Therefore
$L_t-o\le M_t$ whenever $L_t>M_t$, and the proposed quote is always safe
when it is posted.  If $L_t-o>p_t$, the loss against $M_t-p_t$ is
\[
  M_t-L_t+o\le(M_t-L_t)_++o.
\]
If $L_t-o\le p_t$, the player abstains, and
\[
  M_t-p_t
  \le (M_t-L_t)_++(L_t-p_t)
  \le (M_t-L_t)_++o.
\]
Summing proves \eqref{eq:full-upper}.

For $o>0$, set $\gamma=\sqrt{2o/N}$ in
\eqref{eq:random-upper}; this gives
$R_T^{\mathrm{dyn}}\le u+\sqrt{2No}$.  Use whichever of random subtraction
and full subtraction has the smaller displayed guarantee.  Finally, every posted
quote has nonnegative reward and $D_T\le N$, so the dynamic regret of either
policy is also at most $N$.  This proves \eqref{eq:minimax-upper}.  The case
$o=0$ follows directly from full subtraction.
\end{proof}

The absence of a logarithmic term is already visible when the estimates are
exact.  If $L_t=M_t$ and $M_t>p_t$, quote $M_t$; by
\eqref{eq:threshold-equivalence}, it executes later and earns $M_t-p_t$.
If $M_t=p_t$, abstain.  This obtains $D_T$ itself and therefore dominates
every fixed quote.  Its dynamic regret is zero and its fixed-quote regret is
nonpositive.

\section{Matching lower bound}

Define, for $N\ge1$,
\begin{equation}\label{eq:phi}
  \phi_N(o):=\min\{No,\sqrt{No}\},
\end{equation}
and put $\phi_0(0):=0$.  The two terms in \eqref{eq:phi} correspond to two
different ranges: $No$ is the smaller term when $No\le1$, and
$\sqrt{No}$ is the smaller term when $No\ge1$.

\begin{theorem}[Tight dynamic regret]\label{thm:matching}
For every $T\ge2$, $N=T-1$, and $0\le u,o\le N$,
\begin{equation}\label{eq:matching}
  \frac{u+\phi_N(o)}{128}
  \le \Rstar_T^{\mathrm{dyn}}(u,o)
  \le
  \min\left\{
    N,\,
    u+\min\bigl\{No,\sqrt{2No}\bigr\}
  \right\}.
\end{equation}
Consequently,
\begin{equation}\label{eq:rate}
  \Rstar_T^{\mathrm{dyn}}(u,o)
  =\Theta\!\left(
     \min\{N,\,u+\phi_N(o)\}
   \right).
\end{equation}
Moreover, $\Rstar_T^{\mathrm{dyn}}(0,0)=0$ exactly.  When $T=1$, there are
no trades and the value is zero.
\end{theorem}

The upper bound has already been proved.  For the lower bound, we prove the
stronger statement for fixed-quote regret.  The ordering
\eqref{eq:minimax-order} then transfers it immediately to dynamic regret.
We handle underestimation and overestimation separately.

\subsection{Underestimation}

\begin{proposition}\label{prop:under-lower}
For every $N\ge1$ and $0\le u\le N$,
\[
  \Rstar_T^{\mathrm{fix}}(u,0)\ge\frac u4.
\]
\end{proposition}

\begin{proof}
The claim is immediate for $u=0$, so suppose $u>0$.  Put
\[
  h:=\frac uN,
  \qquad
  \ell:=\frac h2.
\]
Consider two price sequences, chosen with equal probability.  Both have
$p_1=\cdots=p_N=0$.  The terminal price is $p_T=\ell$ on the low sequence
and $p_T=h$ on the high sequence.  Reveal $L_t=0$ on every quote round.

On the sequence with terminal price $x\in\{\ell,h\}$, we have $M_t=x$ for
all $t\le N$, so $\cU_T=Nx\le u$ and $\cO_T=0$.  The best fixed quote is
$a=x$ and earns $Nx$; equality with the terminal price executes the orders.

All observations and all fill information available before the terminal
price are identical on the two sequences.  For any quote $a$ selected on a
quote round, its average reward over the two sequences is at most
$\ell$: if $a\le\ell$, it earns at most $\ell$ on both; if
$\ell<a\le h$, it earns $a$ only on the high sequence and hence at most
$h/2=\ell$ on average; and if $a>h$, it earns zero.  The same statement holds
after averaging over a randomized policy.

The player's expected total reward is therefore at most $N\ell=Nh/2$,
whereas the average fixed-quote benchmark is
\[
  \frac{N\ell+Nh}{2}=\frac{3Nh}{4}.
\]
The average regret is at least $Nh/4=u/4$, so at least one of the two
deterministic sequences has expected regret at least $u/4$.
\end{proof}

\subsection{Overestimation}

The following construction will be used in both ranges of $o$.

\newpage
\begin{lemma}[Block construction]\label{lem:block}
Let $n,r\ge1$ be integers satisfying
\[
  n(r+1)+1\le T,
\]
and let $d\in(0,1)$ satisfy $nd\le1/2$.  If $o\ge rd$, then every randomized
policy has expected fixed-quote regret at least
\begin{equation}\label{eq:block-result}
  \frac{dnr}{4}
\end{equation}
on some deterministic sequence with $\cU_T=0$ and $\cO_T\le o$.
\end{lemma}

\begin{proof}
Put $y:=1-d$.  Make $n$ blocks, each consisting of one marker followed by
$r$ zero prices, and add one final marker:
\[
  q_1,\underbrace{0,\ldots,0}_{r},
  q_2,\underbrace{0,\ldots,0}_{r},\ldots,
  q_n,\underbrace{0,\ldots,0}_{r},q_{n+1}.
\]
If this uses fewer than $T$ prices, pad the sequence with copies of the last
marker and reveal the exact estimate on every padding round.  There
$p_t=M_t$, so neither the player nor any fixed quote can earn.

On sequence $H$, every marker equals $1$.  For each $j\in\{1,\ldots,n\}$,
sequence $J_j$ has markers $q_1,\ldots,q_j$ equal to $1$ and all later
markers equal to $y$.  All non-marker prices remain zero.

On $H$, reveal the exact value $L_t=M_t$.  On $J_j$, reveal $L_t=1$ up to and
including all $r$ zero rounds in block $j$, and reveal the exact value $M_t$
afterward.
The estimates before block $j$ are already exact and equal to $1$.  At the
marker starting block $j$, the current price is $p_t=M_t=1$, so its estimate
$L_t=1$ is also exact even though all later markers equal $y$.  Only the $r$
zero rounds in block $j$ have $M_t=y$ and $L_t=1$.  Hence
\begin{equation}\label{eq:block-errors}
  \cU_T=0,
  \qquad
  \cO_T=rd
\end{equation}
on $J_j$, while both errors vanish on $H$.

We next verify the fixed-quote benchmark.  At every marker, including
padding, $p_t=M_t$, so the strict requirement $p_t<a$ makes every genuine
quote unsuccessful.  Fixed quotes therefore earn zero there.  On $H$, quote
$1$ earns one on all $nr$ zero rounds, so
\begin{equation}\label{eq:BH}
  B_H=nr.
\end{equation}
On $J_j$, a quote $a\le y$ can earn on all zero rounds and is maximized at
$a=y$, giving $nry$.  A quote $a>y$ can earn only on the $(j-1)r$ zero
rounds before block $j$, and gives at most $(j-1)r$.  Since $nd\le1/2$
implies $ny\ge n-1$, quote $y$ is optimal and
\begin{equation}\label{eq:BJ}
  B_{J_j}=nry.
\end{equation}

Up to every choice in block $j$, $H$ and $J_j$ reveal the same prices and
estimates.  Quotes from earlier blocks encounter the same markers
$q_2,\ldots,q_j=1$ and hence have identical fill and reward histories.  Quotes
from block $j$ cannot fill before $q_{j+1}$, after all $r$ choices in the
block.  On a zero round, encode abstention by setting $A=0$, and write
$X:=1-A$.
Let
\[
  G_j:=\sum_{t\text{ a zero round in block }j}
         \Pr_H(X_t<d).
\]
The same probabilities apply on $J_j$ during this block.

Write $R_H$ and $R_{J_j}$ for the expected fixed-quote regrets on the
corresponding sequences.  On $H$, every quote placed at a zero price
eventually executes, and quote $1$ is the benchmark.  Therefore
\begin{equation}\label{eq:RH}
  R_H=\sum_{t\text{ a zero round}}\E_H X_t.
\end{equation}
On $J_j$, each of the $(j-1)r$ earlier zero rounds can give the player a
surplus of at most $d$ relative to quote $y$.  During block $j$, the event
$X_t<d$ means $A_t>y$, so the quote fails and loses the benchmark reward
$y$.  If $X_t\ge d$, then either the player abstains or $A_t\le y$; equality
$X_t=d$ means $A_t=y$ and executes by the non-strict rule.  In all these
cases the loss relative to quote $y$ is nonnegative.  All later losses are
also nonnegative.  Thus
\begin{equation}\label{eq:RJ}
  R_{J_j}\ge yG_j-(j-1)rd.
\end{equation}

Let $R$ be the largest expected regret among
$H,J_1,\ldots,J_n$.  Equation \eqref{eq:RJ} gives
\[
  G_j\le\frac{R+(j-1)rd}{y}.
\]
Also, $\E X_t\ge d\Pr(X_t\ge d)$, so \eqref{eq:RH} yields
\begin{align*}
  R
  &\ge d\sum_{j=1}^n(r-G_j)\\
  &\ge dnr-\frac{dn}{y}R
       -\frac{d^2rn(n-1)}{2y}.
\end{align*}
Because $nd\le1/2$, we have $dn/y\le1$, and the last term is at most
$dnr/2$.  It follows that
\[
  R\ge\frac{dnr}{2}-R,
\]
which proves \eqref{eq:block-result}.
\end{proof}

\begin{proposition}\label{prop:over-lower}
For every $N\ge1$ and $0\le o\le N$,
\[
  \Rstar_T^{\mathrm{fix}}(0,o)\ge\frac{\phi_N(o)}{64}.
\]
\end{proposition}

\begin{proof}
The statement is immediate when $o=0$, so assume henceforth that $o>0$.
First suppose $N\ge2$.  One choice
that handles all small and intermediate values is
\[
  r=1,
  \qquad
  n=\left\lfloor\frac N2\right\rfloor,
  \qquad
  d=\min\left\{o,\frac1{2n}\right\}.
\]
Then $rd\le o$, $nd\le1/2$, and Lemma~\ref{lem:block} gives
\begin{equation}\label{eq:small-o-lower}
  R\ge\min\left\{\frac{no}{4},\frac18\right\}
   \ge\min\left\{\frac{No}{12},\frac18\right\},
\end{equation}
where we used $n\ge N/3$.  If $No<64$, the last expression is at least
$\phi_N(o)/64$.

It remains to cover $s:=\sqrt{No}\ge8$.  Choose
\[
  r=\left\lfloor\frac s4\right\rfloor,
  \qquad
  n=\left\lfloor\frac{N}{r+1}\right\rfloor,
  \qquad
  d=\frac1{2n}.
\]
Since $o\le N$, we have $s\le N$.  Moreover,
\[
  r\ge\frac s8,
  \qquad
  r+1\le\frac{3s}{8}\le\frac N2,
  \qquad
  n\ge\frac{N}{2(r+1)}.
\]
The construction fits within the horizon, $nd=1/2$, and
\[
  rd=\frac r{2n}
  \le\frac{r(r+1)}{N}
  \le\frac{3s^2}{32N}
  =\frac{3o}{32}
  \le o.
\]
Lemma~\ref{lem:block} now gives
\[
  R\ge\frac r8\ge\frac s{64}
   =\frac{\phi_N(o)}{64}.
\]

Finally suppose $N=1$ (and still $o>0$).  Compare two sequences with $p_1=0$.  On $H$, let
$p_2=1$; on $J$, let $p_2=1-d$, where $d:=\min\{o,1/2\}$.  Reveal
$L_1=1$ on both.  Thus $\cO_T=0$ on $H$ and $\cO_T=d\le o$ on $J$, with no
underestimation.  Put $X=1-A_1$, again encoding abstention as $A_1=0$, and
let $q=\Pr_H(X<d)$.  The histories before the only choice agree.  The regret
on $H$ is at least $d(1-q)$, whereas on $J$ it is at least $(1-d)q$.
If both quantities were smaller than $d(1-d)$, the first would imply $q>d$
and the second $q<d$, a contradiction.  Their maximum is therefore at least
$d(1-d)\ge o/4$.  Since
$\phi_1(o)=o$ for $o\in[0,1]$, this proves the endpoint.
\end{proof}

\begin{proof}[Proof of Theorem~\ref{thm:matching}]
The hard instances in Proposition~\ref{prop:under-lower} have
$\cO_T=0\le o$, while those in Proposition~\ref{prop:over-lower} have
$\cU_T=0\le u$.  Both families therefore lie in the class defining
$\Rstar_T^{\mathrm{fix}}(u,o)$, and
\[
  \Rstar_T^{\mathrm{fix}}(u,o)
  \ge\max\left\{\frac u4,\frac{\phi_N(o)}{64}\right\}
  \ge\frac{u+\phi_N(o)}{128}.
\]
The minimax ordering \eqref{eq:minimax-order} and the dynamic upper bound in
Theorem~\ref{thm:upper} now prove \eqref{eq:matching}.  Since
$u,\phi_N(o)\le N$, replacing
$u+\phi_N(o)$ by $\min\{N,u+\phi_N(o)\}$ changes it by at most a universal
factor.  This proves \eqref{eq:rate}.

When $u=o=0$, we have $L_t=M_t$.  Quoting $M_t$ when $M_t>p_t$ and
abstaining otherwise earns $D_T$, so its dynamic regret is zero on every
sequence.  A constant price sequence has $D_T=0$ and forces regret zero, so
the worst-case value cannot be negative.  Hence
$\Rstar_T^{\mathrm{dyn}}(0,0)=0$.
\end{proof}

\begin{corollary}[Fixed-quote regret]\label{cor:fixed}
For every $T\ge2$ and $0\le u,o\le T-1$,
\[
  \Rstar_T^{\mathrm{fix}}(u,o)
  =\Theta\!\left(
     \min\{T-1,\,u+\phi_{T-1}(o)\}
   \right).
\]
Moreover, $\Rstar_T^{\mathrm{fix}}(0,0)=0$ exactly.
\end{corollary}

\begin{proof}
Propositions~\ref{prop:under-lower} and~\ref{prop:over-lower} give
\[
  \Rstar_T^{\mathrm{fix}}(u,o)
  \ge \frac{u+\phi_{T-1}(o)}{128}.
\]
The upper bound follows from
$\Rstar_T^{\mathrm{fix}}(u,o)\le\Rstar_T^{\mathrm{dyn}}(u,o)$ and
Theorem~\ref{thm:matching}.  At $(u,o)=(0,0)$, the exact-estimate policy has
fixed-quote regret at most zero on every sequence, while a constant price
sequence forces regret zero.  Thus the minimax value is exactly zero.
\end{proof}

\section{Consequences and interpretation}

\paragraph{No additional estimate.}
Taking $L_t=p_t$ uses only the price already observed at the start of the
round.  Then
\[
  \cU_T=\sum_{t=1}^{N}(M_t-p_t)=\cV_T,
  \qquad
  \cO_T=0.
\]
Consequently, for $0<v\le N$, the minimax value restricted to
$L_t=p_t$ and $\cV_T\le v$ has, for either the dynamic or fixed-quote
benchmark, the tight rate
\begin{equation}\label{eq:variation-corollary}
  \Theta\bigl(\min\{T-1,v\}\bigr).
\end{equation}
The simple dynamic-regret upper bound is obtained by always abstaining:
$D_T=\cV_T$.  The fixed-quote upper bound follows because $B_T\le D_T$.
The terminal-price construction in Proposition~\ref{prop:under-lower} has
$B_T=D_T$ and gives the matching lower bound for both benchmarks.
When $v=0$, both minimax regrets are exactly zero.

\paragraph{One-sided estimates.}
If $p_t\le L_t\le M_t$ and
\[
  W_T:=\sum_{t=1}^{N}(M_t-L_t),
\]
then $\cU_T=W_T$ and $\cO_T=0$.  Quoting $L_t$ whenever $L_t>p_t$ gives
$R_T^{\mathrm{dyn}}\le W_T$, and hence also
$R_T^{\mathrm{fix}}\le W_T$.  For $0<w\le N$, the minimax value
restricted to these one-sided estimates and $W_T\le w$ has, for either
benchmark, the tight rate
\[
  \Theta\bigl(\min\{T-1,w\}\bigr).
\]
When $w=0$, the estimates are exact and both minimax regrets are exactly
zero.

\paragraph{Why the old logarithmic term disappears.}
In the lagged gross-revenue protocol, the estimate of $M_t$ arrived only
after the quote at round $t$, and the player had to leave a multiplicative
margin for the unknown drop in the maximum.  Here $p_t$ and $L_t$ are both
known before the quote.  With an exact estimate, the player can take the
best possible profit on every round, so even a constant logarithmic baseline
would be false.

\paragraph{Why the overestimation term has two ranges.}
If the total overestimation is at most $o$, subtracting $o$ makes every quote
safe and costs at most $No$.  Random subtraction spreads the protection over
the rounds and costs order $\sqrt{No}$.  The block construction shows that
the better of these two costs is unavoidable.  In particular, a square-root
term alone cannot be tight as $o\downarrow0$, because exact estimates have
zero regret.

\begin{remark}[Knowledge of the budget]
The optimized policy uses the announced upper bound $o$, as does the
minimax definition \eqref{eq:minimax}.  This is not merely a proof
convenience.  Already with one quote round, zero dynamic regret on an exact
estimate can force one particular quote, while an arbitrarily close
overestimate makes that same quote fail.  Concretely, take $T=2$, $p_1=0$,
and $L_1=x$.  On the path $p_2=x$, zero regret forces $A_1=x$ almost surely.
On the indistinguishable path $p_2=x-\delta$, that quote fails while both
benchmarks equal $x-\delta$, although $\cO_T=\delta$.  Thus a single policy
that does not know an error scale
cannot retain both exact zero regret at $\cO_T=0$ and a guarantee tending
continuously to zero with $\cO_T$.
\end{remark}

\end{document}
